\section{What the Stability Analysis Means in Plain Language} \label{sec:fsa-v4-intuition}

Section \ref{sec:fsa-v4-stability} contained quite a lot of mathematics. This companion section is for sport-physicians, exercise scientists, coaches and athletes who want the take-home message without the algebra. The picture that emerges is, in our view, surprisingly intuitive — and matches what experienced athletes already know about how their bodies respond to training.

\subsection{The three states of an athlete}

Under the FSA-v4 model, for any constant amount of training load $(\Phi_B, \Phi_S)$ — aerobic and strength stimulus respectively — an athlete's long-term physiology settles into one of three states. Which one depends on the size of the load. Figure \ref{fig:v4_bifurcation_curve} (page \pageref{fig:v4_bifurcation_curve}) is the map.

\paragraph{1. The ``healthy'' state — moderate, sustainable training.}
At low to moderate training loads, the body settles into a state where autonomic amplitude $A$ — a measure of cardiovascular variability and the capacity to mount a stress response — sits at a healthy positive value. Aerobic fitness $B$ and strength $S$ are also present. Crucially, this state is the \emph{only} long-term outcome at these loads: \emph{wherever you start, you end up here}. A poorly-rested, deconditioned, or injured athlete will progressively recover toward this state simply by sticking to the prescribed load. Coaches who say ``trust the process'' are describing this regime: the body has a single attractor and is being pulled toward it.

\paragraph{2. The ``collapsed'' state — chronic overtraining.}
At high sustained loads, the body settles into a state with $A = 0$ — autonomic suppression, blunted heart-rate variability, dampened stress response — even though aerobic and strength capacities $B$ and $S$ may still be high. This is the model's representation of overtraining syndrome / unexplained underperformance / autonomic burnout. At these loads, this state is also the \emph{only} long-term outcome: no matter how rested or motivated the athlete starts out, sustained high training pulls them down to autonomic collapse. The coach's instinct that ``you can't outwork chronic overtraining'' is a statement about this regime.

\paragraph{3. The ``bistable'' state — moderate-to-high training, where it depends on the day.}
Between the healthy and collapsed regimes there is a band of training loads at which \emph{both} states are stable simultaneously. Whether the athlete settles into healthy autonomic function or autonomic collapse depends on \emph{the state they were in} when the training block began. There is a separating threshold (the ``separatrix'' in Figure \ref{fig:v4_basin_phiB}). If the athlete's autonomic amplitude $A$ at the start of the block was above the threshold, they end up healthy. If it was below, they collapse. Same training, two different outcomes.

This is the technical formulation of something every coach has experienced: the \emph{same} training session is constructive on a Tuesday after a good week, and disastrous on a Tuesday after a bad week. In the bistable regime, training tolerance is genuinely state-dependent; it is not noise or imagination.

\subsection{What the stimulus map says you can and cannot do}

The map in Figure \ref{fig:v4_bifurcation_curve} divides the training-load plane into three regions. Going from low load to high load, you cross two boundaries:

\begin{itemize}
  \item The \textbf{transcritical curve} (black solid). Below this curve, the healthy state is the unique attractor — easy life. Above it, the healthy state still exists but \emph{coexists with collapse}.
  \item The \textbf{saddle-node curve} (magenta dashed). Above this, the healthy state has \emph{disappeared entirely} — the body cannot tolerate this load at all, and any starting point ends in collapse.
\end{itemize}

For practical training:
\begin{itemize}
  \item \textbf{Inside the black curve (low load):} you cannot overtrain at this level even if you tried. Useful for rehab, recovery weeks, beginners.
  \item \textbf{Between black and magenta (moderate-to-high load):} you can train hard here \emph{provided you start from a recovered state}. This is where elite training lives. The risk is real but managed.
  \item \textbf{Beyond the magenta curve (very high load):} no matter how recovered you are at the start, sustained loading at this level guarantees collapse. Avoid for sustained periods.
\end{itemize}

\subsection{Sedentary inactivity is not collapse}

A natural question: doesn't $\Phi = 0$ — sedentary detraining — also lead to $A = 0$? In the FSA-v4 model, \textbf{no}: the model has a non-zero baseline autonomic drive ($\mu_0 > 0$ even at zero load), so a fully detrained sedentary individual settles into a low-but-positive $A^*$ rather than zero. Sedentary lifestyles produce \emph{deconditioned but autonomically intact} states ($B \to 0$, $S \to 0$, $A \approx 0.3$). The pathological $A = 0$ state in this model represents \emph{overtraining}, not under-training. (Medical reality is more complicated — long-term sedentary inactivity does eventually impair autonomic function — but the FSA-v4 model is calibrated to the training-stress regime, where the dominant pathology is overload.)

\subsection{What ``optimal training'' means under this model}

The mathematical question is: pick a training programme $\Phi(t)$ to drive the athlete into the healthy basin, deeply, and keep them there with the smallest possible training load. This translates to a specific recipe for the cost function the optimiser should solve. We discuss it in Section \ref{sec:fsa-v4-stability}; the short answer for the practitioner is:

\begin{enumerate}
  \item \textbf{The thing to maximise is the area under the autonomic curve $A(t)$.} Average autonomic amplitude over the training block is the right primary outcome — it tracks both ``in healthy state'' (because $A=0$ in the collapsed basin contributes nothing) and ``deep in healthy state'' (because the deeper the basin, the larger $A^*$).
  \item \textbf{Subtract a small effort penalty.} A cost on the magnitude of training load $\|\Phi(t)\|^2$ ensures that two programmes that produce equal autonomic outcomes are differentiated by which one demanded less from the athlete.
  \item \textbf{If aerobic fitness and strength matter as outcomes in their own right, add weighted $B(t)$ and $S(t)$ terms.} For pure recovery / autonomic-resilience problems, these are not needed; for race-preparation, they are.
  \item \textbf{The hard ``fatigue must not exceed $F_{\max}$'' constraint of FSA-v3 may no longer be necessary.} The variable-dose dynamics of FSA-v4 (the $K_{FB}, K_{FS}$ states) penalise sustained loading naturally — a controller that ignores fatigue progressively raises the fatigue gain and is punished by autonomic collapse without an explicit barrier needing to be written into the cost. A softer penalty on elevated $K$-states is a cleaner v4-native alternative to the hard $F_{\max}$ ceiling.
\end{enumerate}

\subsection{What the controller has to be able to do}

In the bistable regime, the right training plan depends on \emph{the athlete's current physiological state} — the day-to-day $A_t$, $F_t$ values. A plan that ignores those values is a fixed schedule, and a fixed schedule is structurally incapable of the day-by-day adjustments the bistability map requires. This is why Section \ref{sec:open_loop_limitation} concluded that an open-loop optimiser is the wrong tool for FSA-v4. A closed-loop controller — one that continuously updates the plan based on incoming heart-rate-variability, sleep, and performance measurements — is what the model is asking for. The SMC$^2$ control framework (Section \ref{sec:smc2_vs_ot}) is the natural fit because it can simultaneously maintain a probabilistic belief about \emph{where in the basin} the athlete is and \emph{which side of the separatrix} the next training block will land them on.

In short: the basin map gives you an \emph{atlas} of safe and unsafe loads; the bistability tells you tolerance is state-dependent; the variable-dose dynamics tell you tolerance also has memory; and the autonomic amplitude $A$ is the single best summary measure of where in the atlas the athlete currently sits.
